\begin{table}[H]
\centering
\small
\begin{tabular}{@{}lccc@{}}
\toprule
Formule & Neutre & Cation & Anion \\
\midrule
N$_{2}$ & 1 / 3 & 1 / 0 & -- \\
N$_{3}$ & 1 / 13 & 3 / 0 & 3 / 2 \\
N$_{4}$ & 15 / 13 & 1 / 1 & 2 / 1 \quad {\tiny (+1 charge multiple, ex. N$_4^{4-}$)} \\
N$_{5}$ & 2 / 5 & 2 / 17 & 4 / 4 \\
N$_{6}$ & 9 / 8 & 3 / 0 & 2 / 0 \\
N$_{7}$ & 6 / 4 & 2 / 0 & 2 / 0 \\
N$_{8}$ & 26 / 5 & -- & 2 / 0 \\
N$_{9}$ & 2 / 0 & 2 / 0 & 3 / 1 \\
N$_{10}$ & 8 / 9 & 2 / 0 & 1 / 0 \\
N$_{11}$ & 1 / 2 & -- & -- \\
N$_{12}$ & 1 / 5 & -- & -- \\
N$_{13}$ & 0 / 1 & -- & -- \\
N$_{14}$ & 0 / 3 & -- & -- \\
N$_{15}$ & -- & 0 / 1 & -- \\
N$_{16}$ & 0 / 2 & -- & -- \\
N$_{18}$ & 1 / 1 & -- & -- \\
N$_{20}$ & 2 / 1 & -- & -- \\
\bottomrule
\end{tabular}
\caption{Structures confirm\'ees (construites, relax\'ees GFN2-xTB, v\'erifi\'ees en fr\'equence) contre candidats identifi\'es par recherche de citations (format : confirm\'ees / candidats). Les candidats sont compt\'es \`a partir de la formule mentionn\'ee dans le \textbf{titre seul} de l'article -- 78 des 158 candidats (49~\%) n'ont pas de formule identifiable dans leur titre (revues, m\'ethodologie, commentaires) et n'apparaissent donc pas dans cette table malgr\'e une pertinence potentielle. Un compte de candidats n'implique aucune garantie de nouveaut\'e structurale (un article peut tr\`es bien redocumenter une topologie d\'ej\`a connue).}
\label{tab:confirmed-vs-candidates}
\end{table}
